\section{標準 Level Set 法の Godunov 型スキーム（参考）}
\label{app:godunov_ls}
% ═══════════════════════════════════════════════════════════════════════════════

本節は付録~\ref{app:weno5} に収録した WENO5 スキームとは別に，
\textbf{標準レベルセット（符号付き距離関数）}移流に対して使われる Godunov 型 upwinding スキームを説明する．
本稿の主実装（CLS + Dissipative CCD，§\ref{sec:advection_motivation}）では使用しないが，実装の選択肢として参考に記す．

\subsection{Godunov 型スキームの基本形}

符号付き距離関数 $\phi$ の移流方程式 $\phi_t + \bm{u}\cdot\nabla\phi = 0$ に対して，
速度場 $\bm{u}$ の発散フリー条件 $\nabla\cdot\bm{u}=0$ のもとで保存形 $\phi_t + \nabla\cdot(\bm{u}\phi) = 0$ と等価となる．

$x$ 方向フラックスを Godunov 法で評価する：
\begin{equation}
  F_{i+1/2} = u_{i+1/2}^+ \phi_i + u_{i+1/2}^- \phi_{i+1}
\end{equation}
ここで $u^+ = \max(u, 0)$，$u^- = \min(u, 0)$ は速度の上流風分解である．

\subsection{再初期化との組み合わせ}

標準 LS 法では移流後に符号付き距離条件 $|\nabla\phi| = 1$ が崩れるため，
再初期化方程式
\[
  \phi_\tau + S(\phi^0)(|\nabla\phi|-1) = 0
\]
を仮想時間 $\tau$ 方向に数回反復する（$S(\phi^0)$ は初期符号に合わせたスムーズ符号関数）．

CLS 法（本稿主方式）では再初期化が $\psi$ の compact プロファイルを保つよう設計されており，
Godunov 型スキームと組み合わせるよりも質量保存性に優れる（§\ref{sec:reinit}）．

% ═══════════════════════════════════════════════════════════════════════════════
